\chapter{Background and Literature Review}
\label{chap:background}
% ============================================================

\section{Human-Robot Interaction and ROS}

Human-robot interaction (HRI) studies how people and robots communicate, coordinate, and share space. In social and service robotics, the interaction is not only a command interface. The robot must interpret speech, gaze, gesture, people, objects, and task context while remaining physically safe and predictable. Early surveys in HRI and socially interactive robots describe this as a multidisciplinary problem that combines robotics, artificial intelligence, dialogue, perception, and interaction design \cite{fong2003survey,goodrich2007human}. This motivates an architecture in which dialogue, perception, planning, and actuation are not treated as a single monolithic process.

ROS provides a practical software substrate for this separation. A ROS system is built from nodes that communicate through typed topics, services, and actions. Topics are suitable for streams such as perception or status updates, services for request-response queries, and actions for long-running behaviours that need feedback and cancellation. ROS~2 extends this model with design choices aimed at more robust distributed systems, including improved middleware support and deployment considerations for real robots \cite{quigley2009ros,macenski2022ros2}. For this thesis, these boundaries are useful because they allow a language model to remain a high-level planner rather than a direct controller of robot hardware.

ROS4HRI builds on this idea by proposing reusable conventions and interfaces for HRI components in ROS. \textcite{mohamed2020ros4hri} identify the lack of a standard human representation as a barrier to reusable HRI pipelines, particularly when perception modules must represent bodies, faces, voices, groups, and social signals. Their ROS4HRI proposal is relevant here because it treats HRI perception as a set of interoperable streams rather than as robot-specific internal state. A planner should therefore receive compact task-relevant grounding, while perception and skill nodes retain access to richer evidence when they need to execute or validate an action.

This separation also matters for safety. A physical robot should not execute an action only because a language model produced plausible text. The runtime stack should validate that a requested skill exists, that its arguments are admissible, that the plan uses allowed interfaces, and that execution feedback is handled before further actions are dispatched. In this work, \code{nao_orchestrator} is kept as the deterministic execution boundary: it admits planner requests, validates plan steps, dispatches skills, and publishes execution feedback. The language model proposes structured plans, but it does not bypass the ROS execution layer.

\section{NAO Robot Stack}

The physical platform used in this thesis is NAO, a programmable humanoid robot originally developed by Aldebaran Robotics. \textcite{gouaillier2008nao} describe NAO as an autonomous humanoid platform designed for research, education, and RoboCup-style experimentation. Its humanoid form, onboard perception pipeline, speech capabilities, and established use in academic robotics make it suitable for studying HRI behaviours such as spoken interaction, person/object grounding, gaze-oriented actions, navigation-like skills, and feedback-driven task execution. In this thesis, NAO is treated as the embodied execution platform, while the proposed architecture is designed so that the planning layer depends on abstract skills and ROS interfaces rather than on NAO-specific APIs. 


\section{LLMs as High-Level Planners}

Large language models (LLMs) have been studied as high-level planners because they encode procedural and semantic knowledge that can be elicited from natural-language instructions. \textcite{huang2022zeroshot} show that language models can decompose high-level embodied tasks into action-like steps in simulated environments, but also show that raw language-model plans may not map cleanly to admissible actions. This problem is central for robotics: a plan that is linguistically reasonable can still be impossible, unsafe, or unsupported by the current robot.

Several approaches address this limitation by grounding the language model in a restricted action set. SayCan combines LLM task knowledge with affordance scores from pretrained skills so that proposed actions are both semantically appropriate and feasible for the robot \cite{ahn2022saycan}. Code as Policies uses language models to synthesize robot policy code over a provided API, which can express feedback loops and spatial computations, but still depends on a carefully defined set of callable primitives \cite{liang2022codepolicies}. These approaches support the design choice used in this thesis: the LLM should plan over an explicit skill registry, not over arbitrary natural-language actions.

Reasoning and acting methods further show the value of coupling language output with external feedback. ReAct interleaves reasoning-oriented text with actions that query external sources or environments \cite{yao2022react}. In robotics, Inner Monologue uses language feedback from perception, success detectors, and humans to improve long-horizon embodied task execution \cite{huang2022innermonologue}. Reflexion extends the same general idea to language agents that improve across trials through verbalized feedback and memory, without updating model weights \cite{shinn2023reflexion}. These works are relevant because they treat execution feedback as part of the planning process rather than as an afterthought.

This thesis adopts a narrower and more implementable form of closed-loop planning. The planner produces a structured skill plan, \code{nao_orchestrator} validates and executes the plan, skills return structured results, and \code{planner_llm} receives execution feedback. The planner may then request clarification, produce a revised plan, or explain failure. The system does not expose private intermediate reasoning. It exposes only the final structured plan, dialogue acts, and feedback-dependent decisions. This preserves inspectability while avoiding the risks of treating free-form reasoning text as an execution contract.

Prompt-only planning remains limited. A prompt can tell an LLM which skills exist, but it cannot guarantee that the model will always select valid skills, use valid arguments, or remain grounded in the current scene. The stack therefore uses deterministic validation around the LLM. The planner output is JSON, plan steps are checked against allowed skills, recovery policies are constrained, and execution remains bound to the orchestrator. This is the main architectural distinction between using an LLM as a planner and using an LLM as a controller.

\section{Embodied Multimodal Models, VLMs, and VLAs}

Vision-language models (VLMs) and vision-language-action models (VLAs) provide an important context for this work. VLMs connect visual observations with language, while VLAs go further by mapping visual and language inputs to robot actions. BLIP-2 shows how frozen vision encoders and frozen language models can be connected through a lightweight transformer to support vision-language generation \cite{li2023blip2}. LLaVA uses visual instruction tuning to create a multimodal assistant capable of following image-based instructions \cite{liu2023llava}. These systems are not robot controllers, but they show how visual context can be converted into language-compatible representations.

Embodied multimodal models extend this direction into robotics. PaLM-E incorporates visual and continuous sensor inputs into a language model and evaluates the resulting model on embodied tasks including robotic planning and visual question answering \cite{driess2023palme}. RT-2 formulates robotic actions as tokens and co-fine-tunes vision-language models on web-scale data and robot trajectories, giving the model a direct path from visual observations and instructions to action outputs \cite{brohan2023rt2}. OpenVLA makes this family of models more accessible by releasing an open-source VLA trained on a large collection of robot demonstrations \cite{kim2024openvla}.

The present thesis does not attempt to train an end-to-end VLA policy. That would require a different experimental design, robot demonstration data, and low-level control evaluation. The contribution here is instead a modular high-level planning architecture for a ROS4HRI-based NAO stack. VLM and VLA literature is used to motivate future extensions in which richer visual or embodied state could improve the planner's world model. In the current system, visual and KB evidence are converted into compact symbolic context before reaching the planner. This keeps the planner portable across robots, provided that the robot exposes equivalent skills and grounding interfaces.

\section{Symbolic State, Knowledge Bases, and Grounded Context}

Robotic planning requires some representation of the world. Pure perception streams are usually too detailed and unstable for high-level planning, while pure language descriptions risk being incomplete or stale. A symbolic knowledge base (KB) provides one way to bridge this gap. KnowRob represents objects, actions, environments, and task knowledge in a form that supports robot reasoning and querying \cite{tenorth2013knowrob,tenorth2017knowrob}. RoboBrain similarly treats robotic knowledge as a multi-source resource that can support grounding, perception, and planning \cite{saxena2014robobrain}. Other work has explored how interaction with users can help robots acquire knowledge about their environment and ground referential expressions \cite{bastianelli2013knowledge,grassi2021knowledge}.

The current implementation follows the same general principle, but at a smaller scale. Detector-derived observations and KB rows are projected into a compact \code{grounded_context}. This context contains stable entity handles, entity kind, semantic class, visibility, and a bounded set of relevant relations. The planner can then refer to entities such as \code{cup_1} or \code{anonymous_person_ehfbf} without needing raw detector scores, image coordinates, or full RDF-style triples in every prompt.

This distinction is important for HRI. A user may ask ``look at the person'', ``find the cup'', or ``what do you see?''. The robot must connect those phrases to current entities, but the planner does not need to subscribe directly to every perception stream. Skills can still requery raw scene or KB evidence at execution time. This allows the planner prompt to remain compact while preserving access to richer evidence where it is needed. A future World Model Enricher can extend this by maintaining a more explicit evolving scene state, including uncertainty, source recency, and relations between people, objects, and actions. The current system prepares for that extension by keeping grounding as a separate contract rather than embedding it inside the planner logic.

\section{Execution Monitoring, Replanning, and Clarification}

Planning in a physical robot system cannot stop at plan generation. Actions may fail, targets may disappear, the user may change the goal, and the robot may need clarification before continuing. Closed-loop methods in language-agent and embodied-planning literature support this view. Inner Monologue shows that feedback from the environment and humans can improve embodied task completion \cite{huang2022innermonologue}. ReAct shows that language models can combine action with information gathering from external tools or environments \cite{yao2022react}. Reflexion shows that feedback can be stored and reused as language-level experience across trials \cite{shinn2023reflexion}.

The implementation in this thesis uses execution feedback as a runtime contract. A skill result is not only a log message. It informs \code{planner_llm} whether a plan step succeeded, failed, needs user input, or should trigger replanning. The orchestrator never creates a replacement plan locally. If a step fails with a recovery policy such as \code{replan}, \code{nao_orchestrator} blocks unsafe continuation, publishes feedback, and waits for a revised validated plan before dispatching further actions. This maintains a clear separation between execution authority and feedback-driven planning.

Clarification is treated as part of the same loop. If the planner cannot safely select a target, or if execution feedback reports ambiguity, the planner emits a dialogue act rather than inventing a missing fact. The dialogue manager then realizes that act as speech. This makes user interaction part of the execution loop, not a separate fallback. The same mechanism can be used when the user's goal changes during execution, or when the current world state no longer satisfies the preconditions for the next planned action.

\section{Literature Gap}

Existing work demonstrates several important components of language-guided robotics. LLMs can decompose tasks into action-like sequences \cite{huang2022zeroshot}. Affordance-grounded methods can restrict language-model proposals to feasible robot skills \cite{ahn2022saycan}. Multimodal and VLA systems can connect vision, language, and action at larger scale \cite{driess2023palme,brohan2023rt2,kim2024openvla}. Knowledge-base systems show how symbolic state can support grounding and planning \cite{tenorth2013knowrob,saxena2014robobrain}. Execution-feedback approaches show that planning improves when the model can react to environment or human feedback \cite{huang2022innermonologue,yao2022react}.

The gap addressed in this thesis is practical and architectural. We are aware of no existing equivalent ROS4HRI-oriented implementation that combines dialogue routing, compact symbolic grounding, LLM-based high-level planning, deterministic orchestrator validation, skill-level execution, and feedback-driven replanning in a single modular NAO stack. The aim is not to replace ROS execution with a language model. It is to place the language model at the planning and feedback layer while preserving ROS boundaries for perception, skills, validation, and speech realization.

The resulting architecture is positioned between two extremes. It is more flexible than a fixed intent-to-skill mapping because it can produce multi-step plans and react to feedback. It is more controlled than an end-to-end language or VLA policy because it exposes explicit runtime contracts and keeps execution authority inside \code{nao_orchestrator}. This makes it suitable for a master's thesis implementation, while leaving a clear path toward richer world modelling and multimodal grounding in future work.


% ============================================================
